% 全部内容为虚构排版示例，请按实际材料替换。
% 仅供明确要求旧 Word 固定框线范例的场景；新版指南 1.7 使用无框线。
% 在主文件中显式加载 \usepackage[mode=word]{hitstudent-backmatter}。
% 匿名评阅时保持评阅人栏空白。表中主席、委员、秘书的固定行数不可扩充。
\begin{resolutiontext}
\hitwordresolution{%
  \hitwordrow{290.93}{李教授}{教授（是）}{示例大学\\示例学院}{示例学科}
  \hitwordrow{312.65}{示例委员甲}{教授（是）}{示例研究单位}{示例学科}
  \hitwordrow{334.37}{示例委员乙}{教授（是）}{示例大学}{示例学科}
  \hitwordrow{356.09}{示例委员丙}{教授（是）}{示例研究单位}{示例学科}
  \hitwordrow{377.93}{示例委员丁}{教授（是）}{示例大学}{示例学科}
  \hitwordrow{399.65}{示例委员戊}{教授（是）}{示例研究单位}{示例学科}
  \hitwordrow{421.37}{示例委员己}{教授（是）}{示例大学}{示例学科}
  \hitwordrow{443.10}{示例秘书}{副教授（否）}{示例大学}{示例学科}
}{%
本段为虚构排版示例。王同学提交了一篇讨论示例研究问题的学位论文。请将本段替换为实际通过并签署的完整答辩决议，不要沿用本示例的评价内容。

固定表模式先测量每个单元格和决议正文；内容超过固定区域时，编译会报错且不绘制表格。长单位名可以合理换行，较长决议请使用新版无框线原生模式。
}
\end{resolutiontext}
